Many public image datasets are released already augmented: each photograph comes with rotated, flipped, recoloured and filtered copies, and nothing records which files belong together. We show that the standard evaluation of such a release, a random train/test split, measures almost nothing. With $k$ copies per photograph and test fraction $\rho$, a test image has a copy in training with probability $\approx1-\rho^{k-1}$ ($99.2\%$ at $k=4$), so the reported accuracy is an accuracy on seen images. On a controlled release built from CIFAR-10, where fresh photographs give the truth, a random split overstates accuracy by up to \numWeakGapKnn{} points, and with strong edits it \emph{understates} it, because the test images are distorted. We prove that splitting by detected near-duplicate groups leaves test images exposed unless groups are recovered completely, since a group cut into $m$ pieces still leaks with probability $1-\rho^{m-1}$; yet inflation is concentrated on test images whose closest training copy is close in the model's own representation, so incomplete groups can still remove most of it. We introduce AugLeak, a metadata-free detector of hidden copies with an approximately calibrated error rate, built from a filter-robust structure view, transform search, hubness correction and an extreme-value null derived from class labels. It has the highest average precision of all methods we tested on two real releases, and splitting by its groups removes most of the inflation when its views match the edits. An audit of six public collections finds repeat shots filed as independent photographs, test images duplicated in training (some under another label), radiographs filed under two patient identifiers, and a brain-MRI benchmark where about half the test images are seen slices. Code and an auditing tool are released.
